\section{Introduction}

Absence seizures are brief, generalized epileptic events that punctuate the conscious experience of patients with idiopathic generalized epilepsies and that interrupt ongoing behavior without overt convulsive motor signs. Electroencephalographically, they are defined by bilateral, synchronous, and stereotyped spike-and-wave discharges (SWDs), with a frequency of 2--5 Hz in humans and 7--12 Hz in rodent genetic models, that emerge abruptly from an otherwise unremarkable background and persist for a few seconds to several minutes \citep{szaflarski2010,crunelli2002}. Although the events are, by definition, self-limited, their repetition across a day and across decades of a patient's life imposes a substantial cognitive and social cost. A central and still incompletely understood feature of absence seizures is the transient lapse of sensory awareness that accompanies the discharge: patients lose contact with their surroundings, fail to respond to questions, and later have no recollection of stimuli presented during the event \citep{blumenfeld2012,ferri2025}. Understanding how a seizure can so abruptly and reversibly disable conscious perception, while leaving the basic architecture of thalamocortical sensory loops structurally intact, is a problem that sits at the intersection of systems neuroscience and clinical neurology.

The Genetic Absence Epilepsy Rat from Strasbourg (GAERS) recapitulates the core electrographic, behavioral, and pharmacological features of human typical absence epilepsy and has become one of the most extensively characterized animal models of the disorder \citep{marescaux1992,danober1998}. In GAERS, SWDs arise spontaneously from an initiation zone in the perioral region of the primary somatosensory cortex and propagate through thalamocortical loops that include the ventral posteromedial nucleus, the reticular thalamic nucleus, and widespread cortical territories \citep{polack2007,meeren2002}. Intracellular and multi-unit recordings have demonstrated that sensory-evoked firing is not uniformly abolished during SWDs; rather, responses survive in subcortical nuclei and in early cortical layers but fail to be integrated in the way that typically produces a behavioral report \citep{chipaux2013}. These observations motivate a whole-brain, spatially unbiased assessment of how SWDs reshape the routing of sensory information, a question that is inherently well-suited to functional magnetic resonance imaging.

Progress on this front has been limited by three interlocking technical barriers. First, conventional echo-planar imaging (EPI) produces gradient noise in excess of 100 dB, which is aversive to unanesthetized rodents and provokes stress-related autonomic and motion artifacts \citep{paasonen2018}. Second, the anesthetic regimens that mitigate this problem also suppress or deform SWDs, so that imaging under anesthesia measures a qualitatively different phenomenon from the one that clinicians encounter in awake patients \citep{mishra2011}. Third, even in successful awake rodent fMRI preparations, framewise motion and physiological noise typically dominate the small sensory-evoked signals that one hopes to measure. Recent advances in silent or quiet acquisition sequences, including multi-band SWIFT and zero echo time (ZTE) imaging, have begun to address the first barrier by reducing acoustic output to levels that awake rats tolerate after habituation \citep{paasonen2020,wiesinger2022,lehto2017,stenroos2018}. These approaches open a path to mapping spontaneous SWDs and their interaction with external sensory inputs in a fully awake brain.

Here we combine awake concurrent EEG-fMRI in GAERS rats with a silent 3D ZTE sequence to ask how absence seizures modulate sensory-evoked hemodynamic responses across the whole brain. We delivered calibrated visual and whisker stimulation during spontaneously occurring ictal and interictal periods and compared the resulting statistical parametric maps and region-of-interest response amplitudes. Our experimental design yielded 14 visual and 14 whisker stimulation sessions across 11 animals, with framewise head translation held to a few tenths of a micrometer and motion occurrences less than half those reported in previous awake rat fMRI studies using EPI. The approach allowed us to resolve interictal responses in canonical visual and somatosensory pathways and to track, on a voxel-by-voxel and region-by-region basis, how these responses are transformed when a seizure is ongoing at the moment of stimulation.

We make three contributions. First, we provide the first whole-brain hemodynamic map of sensory responsiveness during awake spontaneous absence seizures in a rat model, showing that SWDs are accompanied by a regionally specific collapse of positive BOLD responses in primary sensory, associative, and medial frontal cortices. Second, we quantify the transformation of the evoked response: extreme beta-values in the visual cortex region of interest fall from $2.8 \pm 1.7$ during baseline to $-6.0 \pm 2.0$ during seizures ($p < 0.001$), and analogous changes are seen in the barrel cortex, where responses move from $4.1 \pm 1.9$ to $-9.0 \pm 1.9$ ($p < 0.001$). Third, we show that the same stimulation, when it happens to terminate a seizure, elicits a restored and even enhanced positive response ($8.1 \pm 7.0$ in the visual cortex and $4.8 \pm 2.9$ in the barrel cortex), which argues that the cortex retains the capacity to process the stimulus and that the ictal state specifically gates access to that capacity. Taken together, the findings support a model in which the `wave' phase of the SWD acts as a recurrent, thalamocortically enforced refractory window that intermittently denies sensory input its usual cortical amplification.

\section{Related Work}

Early fMRI studies of absence epilepsy were performed almost exclusively in humans with simultaneous scalp EEG, using EPI acquisition during spontaneous SWDs recorded opportunistically in the scanner \citep{gotman2005,moeller2008}. These reports established a robust phenomenology: widespread cortical deactivations, particularly in default-mode network nodes, accompanied by thalamic activations. High-density EEG-fMRI work by \citet{tenney2014} separated low- and high-frequency SWD subnetworks and suggested that distinct oscillatory components recruit different brain regions. However, human studies are constrained by the opportunistic nature of SWD occurrence, the difficulty of delivering precisely timed sensory stimuli during ictal windows, and the limited ability to control for fluctuations in arousal. They also provide no access to single-unit or layer-resolved intracortical signals that would connect the hemodynamic observations to the generative circuit.

Rodent fMRI of SWDs, in turn, has historically been performed under anesthesia, most commonly in WAG/Rij rats. \citet{mishra2011} provided the paradigmatic study, demonstrating that ictal BOLD changes in cortex and thalamus, while consistent with human data, dissociate from neuronal firing in ways that complicate a straightforward neurovascular interpretation. These observations dovetail with the electrophysiological work of \citet{polack2007} and \citet{meeren2002} on the cortical focus hypothesis and with intracellular recordings from \citet{chipaux2013} that show preserved early sensory responses during SWDs. None of these prior studies, however, combined awake imaging, spontaneous SWDs, and precisely timed multimodal sensory stimulation in a single preparation.

Finally, our work sits within a broader methodological push toward silent or low-noise fMRI acquisition in awake rodents. \citet{paasonen2018,paasonen2020} used multi-band SWIFT to show that sensory and default-mode networks can be measured in awake rats with motion levels well below those reported for EPI; \citet{lehto2017} demonstrated MB-SWIFT during deep brain stimulation. Zero echo time imaging, reviewed in \citet{wiesinger2022}, represents a complementary silent strategy in which the readout begins during the ramp of the excitation pulse, effectively eliminating acoustic gradient switching. Our study adapts this family of methods to the specific demands of imaging spontaneous pathological rhythms in awake animals, building on the awake rat restraint and habituation procedures introduced by \citet{stenroos2018} and the motion correction framework of \citet{power2012}.
